%% notes-tex/031-unified-representation/sections/9-model.tex

\section{Where the channels enter the cell graph transformer\secstatus{todo}}
\label{sec:model}

The representation has to land in a model that already exists, and the model's own structure
decides where each channel can go. This section states what the cell graph transformer does
today, read from \file{torchcell/models/equivariant\_cell\_graph\_transformer.py}, and then
places the genotype dose, the environment context and the dataset token in it. Anything not
yet run is labeled as a hypothesis.

\notesec{The stack as it stands}
The encoder holds one token per gene, $N = 6{,}607$, plus a CLS token, $H^{(0)} \in
\mathbb{R}^{(N+1) \times d_h}$, from a learnable table or from precomputed sequence
embeddings. $L$ transformer layers follow, and in each layer every attention head is pulled
toward one of the nine gene-gene graphs by a divergence penalty
$\mathrm{KL}(A_g \,\Vert\, \alpha^{(\ell,k)})$ of weight $\lambda$, or restricted to that
graph's edges by a hard mask. The perturbation then enters through a Type I operator,
representation to representation: the perturbed genes' current embeddings are the keys and
values of a cross-attention over every gene token,
\[
  H_{\mathrm{pert},i} \;=\; h_i \;+\; \sum_{j \in S} \alpha_{ij}\, W_V h_j ,
\]
so the output keeps its $[\,\text{batch}, N, d_h\,]$ shape and is equivariant over genes.
The operator is neither a mask nor a divergence term. It is soft by construction, and every
perturbation in it is a set membership, $j \in S$: the model knows which genes are perturbed
and nothing about how much. Type II readouts then map $H_{\mathrm{pert}}$ to outputs, and
the existing ones are the instruments the earlier experiments built: a pooled perturbation
head for gene interaction (experiment 010), a global head for fitness, a per-gene head for
expression and proteome that keeps the per-node structure and can emit distributional
parameters (019), and a per-metabolite head that pools gene embeddings over the
gene-reaction-metabolite incidence of Yeast9 (026 and 027), with a masked multitask loss
that scores each head only on the records carrying its phenotype. Experiment 030 added a
per-entry dataset token concatenated to the head input at the readout.

\notesec{Genotype: the dose scales the operator}
The operator's set membership generalizes to a scaled membership without changing its form.
With $m_j = 1 - d_j$ the perturbation magnitude at gene $j$, one for a deletion and one half
for a heterozygous deletion in a diploid,
\[
  H_{\mathrm{pert},i} \;=\; h_i \;+\; \sum_{j \in S} \gamma(m_j)\, \alpha_{ij}\, W_V h_j ,
  \qquad \gamma(1) = 1 .
\]
The constraint $\gamma(1) = 1$ is what makes this the smaller of the two changes: every
record in every current dataset is a deletion with $m_j = 1$, so every trained model and
every earlier result is reproduced exactly, and the only new behavior is at $m_j =
\tfrac{1}{2}$. Whether $\gamma$ is a fixed identity, $\gamma(m) = m$, or a learned scalar
map is an arm, not a design choice; the 648{,}977 Hoepfner cells measured in both arms are
the data that decide it, and they are the only such data. \textit{Hypothesis (untested):} a
heterozygote attenuates the deletion response rather than producing a different one, which
is what a scalar $\gamma$ can express and a separate feature would have to learn. The same
$d_i$ replaces the presence bit the flux layer uses to zero a deleted gene's reactions, so
the metabolic branch stops asserting a full knockout of a gene with a working copy.

\notesec{Environment: nothing reaches the tensor layer today}
The dataset-to-tensor processor reads exactly one field off every perturbation, the
systematic gene name, and no environment field at all: no compound, no medium, no dose
appears anywhere in \file{torchcell/data/graph\_processor.py}. Every chemogenomic record
therefore looks to the current model like a deletion with no condition. The environment
context $z_E$ of Section~\ref{sec:representation}, the pooled compound embeddings, the
pooled medium embeddings, the dose-basis one-hot and the masked log molar value, has two
possible entry points and they are the first two arms.

The first is the readout: $z_E$ is concatenated to the Type II head input beside the
dataset token, exactly as 030 concatenates the token. This is the smallest change and it
leaves the encoder and the operator untouched, so the gene representation is condition-free
and the head does all the conditioning. The second is the operator: the compound embeddings
are projected to $d_h$ and appended as extra key-value rows of the cross-attention, so every
gene token attends to the compound as it attends to the perturbed genes, and the compound
shapes the perturbed representation itself. \textit{Hypothesis (untested):} the second arm is
the one that can represent a gene-by-compound interaction, which the readout arm can only
approximate through the pooled $z_S$, and it is where the 0.31-against-0.09 gap of
Section~\ref{sec:established} says the signal lives; the first arm is the control that says
how much the readout alone recovers.

\notesec{One decoder over four response scales}
The four kept datasets report a response on three scales, a log ratio, a z score and a
sensitivity score, from two assays (Table~\ref{tab:physical}). Three decoder designs are
available. Four heads, one per dataset, under the existing masked multitask loss: this is
what the code does today for different phenotype types, and it shares nothing at the output.
One head with a per-dataset output calibration, a learned affine map per source: it shares
the head and lets each source have its scale. Or one head with the dataset token at the
readout, the 030 mechanism: the token is projected and concatenated to the head input, so
the head learns each source's scale and offset from the same weights.

The third is the one taken here, for a measured reason and a structural one. The measured
reason is the 030 smoke test on GilaHyper job 2867 (W\&B run \texttt{5xn8w427}): a synthetic
per-source offset of 0.3 injected into twin copies of four Kuzmin sources was recovered as a
prediction difference of 0.285 with standard deviation 0.053, so the token does carry a
source's residual scale. The structural reason is Section~\ref{sec:datasets}: every physical
field is constant within a dataset, so the token is also the only channel those fields need,
and four heads would have to learn them four times. Each dataset's response is standardized
to zero mean and unit variance within the training split before the loss, which the 030
label-normalization transform already does per source, so the three scales meet the loss on
one footing and the token explains what standardization cannot. \textit{Hypothesis
(untested):} whether pooling under this decoder helps Vanacloig at all, against training on
Vanacloig alone, is the first question the model phase answers, and the compound-cold split
of Section~\ref{sec:established} with the reliability ceiling beside every compound is how it
is scored.

\notesec{Masking against the divergence prior}
The 025 graph-regularization sweep ran the nine graphs into layer one under the divergence
prior at seven weights and under the hard mask, three seeds each, on the 376{,}732 triples of
experiment 010 (\texttt{notes-tex/025-graph-reg-sweep}). Read at a fixed epoch 29 the mask
sat with no penalty, $0.421 \pm 0.003$ against $0.425 \pm 0.000$ held-out Pearson, and the
prior rose with $\lambda$ without a peak inside the ladder, to $0.447 \pm 0.001$ at $\lambda
= 1$. Read at each run's best epoch, all arms lay within 0.010 of one another, and what the
prior bought at the fixed reading was mostly resistance to the decline the unpenalized model
showed after epoch 16. That is the evidence behind the suspicion that the mask is imposed too
hard and too early rather than being wrong: a mask that starts open and closes over training,
or a prior whose $\lambda$ anneals upward, has not been run. \textit{Hypothesis (untested):}
an annealed mask reaches the prior's fixed-epoch score. It is a cheap arm to add to the first
round because it changes the schedule and nothing else. The default for this phase is the
prior at $\lambda = 1$, because it is the measured best at the reading that matters for a
model trained to a budget rather than to a peak.

\notesec{What the other instruments contribute}
The chemogenomic head is one more Type II instrument on the same $H_{\mathrm{pert}}$ that
the expression, metabolite and gene-interaction heads read, so a model trained here can be
trained jointly with those datasets under the masked loss, and the environment context is
new to all of them: the 019 expression data and the 010 triples were measured in one
condition each and have never had a condition channel. \textit{Hypothesis (untested):} joint
training with the deletion-fitness and expression data supplies the gene prior that
Section~\ref{sec:established} measured as the transferable part of the chemogenomic signal.
That is a second-round question, after the single-task result exists.
